\begin{lemma}[determinant characters and the Hodge line]
\label{lem:hodge-line-audit}
\Verified\ Let
\[
 \mathfrak K=f_*\omega_{X/B}.
\]
Then \(\mathfrak K\) is invertible and
\[
 \mathfrak K\simeq\mathcal O_B(1).
\tag{9.1}
\]
More precisely,
\[
 \det R_{g_1}=-\frac1\tau,\qquad
 \det R_{g_2}=\frac1\tau,\qquad
 \det R_{g_0}=1,
\tag{9.2}
\]
and at the fixed points
\[
 \det R_{g_j}(z_j)=\zeta_j^{a_j},
 \qquad (a_1,a_2)=(2,1).
\tag{9.3}
\]
The form \(e\) generates \(\mathfrak K\) at the cusp, while a unit times
\(s_j^{-a_j}e\) generates it at \(p_j\).
\end{lemma}

\begin{proof}
Equation (2.3d) in \cref{prop:period-input-audit} gives
\[
 R_g(z)\Pi(z)\lambda=\Pi(gz)A_g\lambda.
\]
Using
\[
 A_1\widehat w=\widehat u-\widehat w,\quad
 A_1\widehat\delta=\widehat\delta,\qquad
 A_2\widehat w=-\widehat u+\widehat\delta,\quad
 A_2\widehat\delta=\widehat\delta,
\]
and writing \(\tau^\sharp=\tau(g_jz)\),
\(\mu^\sharp=\mu(g_jz)\), direct multiplication gives
\[
 R_{g_1}=
 \begin{pmatrix}\tau^\sharp-1&0\\ \mu^\sharp&1\end{pmatrix},
 \qquad
 R_{g_2}=
 \begin{pmatrix}-\tau^\sharp&0\\1-\mu^\sharp&1\end{pmatrix}.
\]
The transformation laws
\(\tau\circ g_1=(\tau-1)/\tau\) and
\(\tau\circ g_2=-1/\tau\) prove (9.2) at \(g_1,g_2\); the cusp formulas
give \(R_{g_0}=I\).  Since \(\widehat\delta\) is fixed, one eigenvalue is
\(1\).  At \(\tau(z_1)=\rho=e^{\pi i/3}\), the other is
\(\rho-1=e^{2\pi i/3}=\zeta_1^2\); at \(\tau(z_2)=i\), it is
\(-i=\zeta_2\).  This proves (9.3).

We determine the extension over a finite filling without suppressing the
ramification.  On \(J'_j\), the nowhere-zero form \(ds_j\wedge e\)
transforms by
\[
 (\widetilde g_j^{\log})^*(ds_j\wedge e)
 =\chi_j(s_j)\,ds_j\wedge e,
 \qquad \chi_j=\zeta_j\det R_{g_j}.
\tag{9.4}
\]
Iteration \(m_j\) times shows
\(\prod_{r=0}^{m_j-1}\chi_j(\zeta_j^rs)=1\).
Write \(\chi_j(s)=\chi_j(0)e^{\psi(s)}\), with \(\psi(0)=0\).
The product identity says that the power series of \(\psi\) has no terms
whose degree is divisible by \(m_j\).  Hence
\[
 \varphi(s)-\varphi(\zeta_js)=\psi(s)
\]
has a holomorphic solution, obtained coefficientwise by division by
\(1-\zeta_j^n\).  Therefore
\[
 s^{k_j}e^{\varphi(s)}ds\wedge e,
 \qquad k_j=m_j-1-a_j,
\tag{9.5}
\]
is invariant: \(\chi_j(0)=\zeta_j^{1+a_j}\), so invariance is equivalent
to \(k_j+1+a_j\equiv0\pmod {m_j}\).  We now distinguish the descended
canonical section from the generator of its direct image.  Let
\[
 q_j:J'_j\longrightarrow N_j,\qquad
 \widetilde f_j:J'_j\longrightarrow\Delta_j,\qquad
 \varpi_j:\Delta_j\longrightarrow D_j
\]
denote the finite quotient, the torus-family projection, and the cyclic base
quotient, respectively, so that
\[
 f_j\circ q_j=\varpi_j\circ\widetilde f_j,
 \qquad t_j=s_j^{m_j}.
\]
Put \(\mathcal A_j=\widetilde f_j^{-1}(0)=\{s_j=0\}\) and
\(G_j=\langle\widetilde g_j^{\log}\rangle\).  The action is free by
\cref{prop:candidate-finite-fillings}; hence \(q_j\) is finite \'etale
and \(q_j^{-1}(S_j)=\mathcal A_j\).

The invariant form in (9.5), now denoted
\[
 \sigma_j=s_j^{k_j}e^{\varphi(s_j)}ds_j\wedge e,
 \qquad k_j=m_j-1-a_j,
\tag{9.5a}
\]
descends uniquely to
\(\overline\sigma_j\in H^0(N_j,K_{N_j})\) with
\(q_j^*\overline\sigma_j=\sigma_j\).  Since \(e^{\varphi(s_j)}\) is a
unit and \(ds_j\wedge e\) is nowhere vanishing,
\[
 \operatorname{div}_{J'_j}(\sigma_j)=k_j\mathcal A_j.
\]
Because \(q_j\) is \'etale,
\[
 q_j^*\operatorname{div}_{N_j}(\overline\sigma_j)
 =\operatorname{div}_{J'_j}(q_j^*\overline\sigma_j)
 =k_j\mathcal A_j=q_j^*(k_jS_j).
\]
Faithfulness of finite surjective pullback on Cartier divisors therefore gives
\[
 \operatorname{div}_{N_j}(\overline\sigma_j)=k_jS_j.
\tag{9.5b}
\]
Thus (9.5a) is a canonical section with that divisor, not a nowhere-vanishing
frame when \(k_j>0\); explicitly, \(k_1=0\) and \(k_2=2\).

For the direct-image generator, let
\[
 \vartheta_j=(ds_j\wedge e)\otimes(dt_j)^{-1}
\tag{9.5c}
\]
be the nowhere-vanishing frame of
\[
 q_j^*\omega_{N_j/D_j}
 =K_{J'_j}\otimes(f_j\circ q_j)^*K_{D_j}^{-1}.
\]
Here \((dt_j)^{-1}\) is the pulled-back frame of the inverse canonical line
of \(D_j\), not division by the differential of \(s_j^{m_j}\).  Put
\(\vartheta'_j=e^{\varphi(s_j)}\vartheta_j\).  The defining equation for
\(\varphi\), together with
\(\chi_j(s)=\chi_j(0)e^{\psi(s)}\) and
\(\chi_j(0)=\zeta_j^{1+a_j}\), gives
\[
 (\widetilde g_j^{\log})^*\vartheta'_j
 =\zeta_j^{1+a_j}\vartheta'_j.
\tag{9.5d}
\]
Every holomorphic coefficient of a relative top form is pulled back from
\(\Delta_j\), because the fibres of \(\widetilde f_j\) are connected compact
complex tori.  Hence
\[
 \widetilde f_{j*}q_j^*\omega_{N_j/D_j}
 =\mathcal O_{\Delta_j}\vartheta'_j.
\]
For \(F(s_j)\in\mathcal O_{\Delta_j}\), the section
\(F(s_j)\vartheta'_j\) is \(G_j\)-invariant exactly when
\[
 F(\zeta_js_j)\zeta_j^{1+a_j}=F(s_j).
\tag{9.5e}
\]
Writing \(F(s_j)=\sum_{n\ge0}c_ns_j^n\), this says that \(c_n=0\) unless
\[
 n+1+a_j\equiv0\pmod {m_j},
 \qquad\text{equivalently}\qquad n\equiv k_j\pmod {m_j}.
\]
Consequently,
\[
 \left((\varpi_j)_*\widetilde f_{j*}
 q_j^*\omega_{N_j/D_j}\right)^{G_j}
 =\mathcal O_{D_j}\cdot s_j^{k_j}\vartheta'_j,
\]
and finite descent gives the precise invariant-module identity
\[
 f_{j*}\omega_{N_j/D_j}
 =\mathcal O_{D_j}\cdot s_j^{k_j}\vartheta'_j.
\tag{9.5f}
\]
Finally, over \(\Delta_j^*\), the differential identity
\[
 dt_j=m_js_j^{m_j-1}ds_j
\]
identifies this generator with
\[
 s_j^{k_j}\vartheta'_j
 =\frac{e^{\varphi(s_j)}}{m_j}
 s_j^{k_j-(m_j-1)}e
 =\frac{e^{\varphi(s_j)}}{m_j}s_j^{-a_j}e.
\tag{9.5g}
\]
Thus a unit times \(s_j^{-a_j}e\) generates
\(f_{j*}\omega_{N_j/D_j}=(f_*\omega_{X/B})|_{D_j}\), whereas the
canonical section (9.5a) itself has divisor \(k_jS_j\).

At the cusp, in every smooth toric chart the unimodular monomial change of
coordinates turns
\[
 \Omega_0=\frac{dx_1}{x_1}\wedge
          \frac{dx_2}{x_2}\wedge dt_c
\]
into a unit times the ordinary coordinate volume form.  Thus it generates
the canonical line.  It is invariant under translations and under the
unimodular shears defining the quotient.  Under the punctured-overlap map,
\(\Omega_0=(2\pi i)^2dt_c\wedge e\), so \(e\) generates
\(\mathfrak K\) at \(p_0\).

It remains to compute the degree.  The modular-form calculation already
used in the proof of \cref{prop:period-input-audit} gives
\[
 \Theta=
 \left(\frac{E_4^2E_6^{1/2}}{\Delta_{\rm mod}}\right)\!\circ\tau.
\]
The square root exists because the zeros of \(E_6\circ\tau\) have even
order.  The weights \(8,3,12\) give the same automorphy factor as \(e\);
the sign at \(z_2\) is fixed by its linearizing coordinate.  The orders are
\[
 \operatorname{ord}_{z_1}\Theta=2=a_1,\qquad
 \operatorname{ord}_{z_2}\Theta=1=a_2,\qquad
 \Theta=t_c^{-1}\cdot(\text{unit})
\tag{9.6}
\]
at the cusp.  Consequently \(\mathsf e=\Theta^{-1}e\) is single-valued.
At \(p_j\) it is a unit times \(s_j^{-a_j}e\); at \(p_0\) it is a unit
times \(t_ce\).  Hence it is a global section of \(\mathfrak K\) with the
single simple zero \(p_0\).  This proves
\(\mathfrak K\simeq\mathcal O_B(p_0)\simeq\mathcal O_B(1)\).
Finally, \(\mathfrak K\) is torsion free of generic rank one on the smooth
curve \(B\), hence is invertible.
\end{proof}
